\SourcePage{129}
\SGAsection{1}{扩张与上同调}
\textbf{1.1。}设 $\mathcal C$ 是对\Term{fiber-product}{纤维积}封闭的\Term{category}{范畴}。\footnote{编者注（35）：我们补加了 $\mathcal C$ 对\Term{fiber-product}{纤维积}封闭这一假设（应用中，$\mathcal C$ 是\Term{scheme}{概形}\Term{category}{范畴} $(\mathrm{Sch})$ 或\Term{functor}{函子} $(\mathrm{Sch})^\circ\to(\mathrm{Ens})$ 的\Term{category}{范畴}，均满足此条件）。若略去此假设，以下须假设\Term{fiber-product}{纤维积} $G\times_S\cdots\times_S G$ 都\Term{representable}{可表}。}设 $S$ 为 $\mathcal C$ 的一个\Term{object}{对象}，$G$ 为一个（\Term{representable}{可表}的）$S$-\Term{group}{群}，$F$ 为一个带有 $G$ \Term{action}{作用}的交换 $S$-\Term{group-functor}{群函子}。第 I 讲 5.1 定义了\Term{cohomology}{上同调群} $H^n(G,F)$。回顾其定义：它们是记为 $C^*(G,F)$ 的\Term{complex}{复形}的\Term{homology}{同调群}，其中，令 $(G/S)^n=G\times_S\cdots\times_S G$（$n$ 个因子），则
\[
 C^n(G,F)=\operatorname{Hom}_S((G/S)^n,F).
\]
由于 $G$ 因而\footnote{参见编者注（35）。}各 $(G/S)^n$ 都\Term{representable}{可表}，也有
\[
 C^n(G,F)=F((G/S)^n).
\]
由此及\Term{boundary-operator}{边缘算子}的定义可见，\Term{complex}{复形} $C^*(G,F)$ 只依赖 $F$ 在 $\mathcal C/S$ 的一个\Term{full-subcategory}{全子范畴}上的\Term{restriction}{限制}，该\Term{category}{范畴}的\Term{object}{对象}是 $G$ 在 $S$ 上的\Term{cartesian-power}{笛卡儿幂}。因此有：
\begin{SGAresult}{\Term{lemma}{引理}}{1.1.1}
设 $\mathcal C$ 是对\Term{fiber-product}{纤维积}封闭的\Term{category}{范畴}\footnote{参见编者注（35）。}，$S$ 为 $\mathcal C$ 的\Term{object}{对象}，$G$ 为\Term{representable}{可表}的 $S$-\Term{group}{群}。以 $\mathcal C(G)$ 表示 $\mathcal C/S$ 中以 $G$ 在 $S$ 上的\Term{cartesian-power}{笛卡儿幂}为\Term{object}{对象}的\Term{full-subcategory}{全子范畴}。设 $F,F'$ 是带有 $G$ \Term{action}{作用}的两个交换 $S$-\Term{group-functor}{群函子}。若 $F,F'$ 在 $\mathcal C(G)$ 上的\Term{restriction}{限制}相同，则有\Term{canonical}{典范}\Term{isomorphism}{同构}
\[
 H^*(G,F)\xrightarrow{\sim}H^*(G,F').
\]
\end{SGAresult}
\textbf{1.1.2。\Term{cohomology}{上同调}与\Term{restriction-of-scalars}{标量限制}。}\footnote{编者注（36）：本段标题为新增。}再给出一个比较结果。现在设 $T\to S$ 是 $\mathcal C$ 的一个\Term{morphism}{态射}。若 $F$ 是交换 $T$-\Term{group-functor}{群函子}，则由“\Term{restriction-of-scalars}{标量限制}”（见第 II 讲第 1 节）得到的\Term{functor}{函子}
\[
 F_1=\prod_{T/S}F
\]
是交换 $S$-\Term{group-functor}{群函子}，而且有 $S$-\Term{group-functor}{群函子}的\Term{morphism}{态射}
\[
 u:\prod_{T/S}\underline{\operatorname{Aut}}_{T\text{-gr.}}(F)
 \longrightarrow\underline{\operatorname{Aut}}_{S\text{-gr.}}(F_1).
\]\footnote{编者注（37）：事实上，设 $S'\to S$，$\alpha\in\operatorname{Aut}_{(T\times_S S')\text{-gr.}}(F_{T\times_S S'})$。对每个 $U\to T\times_S S'$，以 $\alpha_U\in\operatorname{Aut}_{\mathrm{gr.}}(F_{T\times_S S'}(U))$ 表示由 $\alpha$ 定义的元素；则 $u(S')(\alpha)$ 是 $\underline{\operatorname{Aut}}_{S\text{-gr.}}(F_1)(S')=\operatorname{Aut}_{S'\text{-gr.}}(\prod_{T\times_S S'/S'}F_{T\times_S S'})$ 中的元素 $\beta$，其\Term{characteristic}{特征}是对每个 $S''\to S'$ 均有 $\beta_{S''}=\alpha_{T\times_S S''}$。}
现在设 $G$ 是 $S$-\Term{group-functor}{群函子}，且有
\[
 G_T\longrightarrow\underline{\operatorname{Aut}}_{T\text{-gr.}}(F),
\]
\SourcePage{130}
即 $G_T$ 对 $F$ 的\Term{action}{作用}。由\Term{functor}{函子} $\prod_{T/S}$ 的定义，得到 $S$-\Term{group-functor}{群函子}的\Term{morphism}{态射}
\[
 G\longrightarrow\prod_{T/S}\underline{\operatorname{Aut}}_{T\text{-gr.}}(F),
\]
再与 $u$ 复合，便得到 $G$ 对 $F_1=\prod_{T/S}F$ 的\Term{action}{作用}。\footnote{编者注（38）：具体地说，对每个 $S'\to S$，$G(S')$ 对 $F_1(S')=F(T\times_S S')$ 的\Term{action}{作用}，由 $G_T(T\times_S S')=G(T\times_S S')$ 对 $F(T\times_S S')$ 的\Term{action}{作用}，以及对应于\Term{projection}{投影} $T\times_S S'\to S'$ 的\Term{group-morphism}{群态射} $G(S')\to G(T\times_S S')$ 给出。也可以说，$G_T$ 对 $F$ 的\Term{action}{作用}由一个\Term{morphism}{态射} $G_T\times_T F\to F$ 给出；施以\Term{functor}{函子} $\prod_{T/S}$，并以 $G_1$ 表示 $S$-\Term{group}{群} $\prod_{T/S}G_T$，由第 II 讲 1.2 得到\Term{morphism}{态射} $G_1\times_S F_1\to F_1$，它定义 $G_1$ 对 $F$ 的\Term{action}{作用}。于是，$G$ 的\Term{action}{作用}经由自然\Term{morphism}{态射} $G\to G_1$ 得到，后者由\Term{adjunction}{伴随}对应于 $\operatorname{id}_G$。}
\begin{SGAresult}{\Term{lemma}{引理}}{1.1.2}
在上述条件下，有\Term{canonical}{典范}\Term{isomorphism}{同构}
\[
 H^*\left(G,\prod_{T/S}F\right)\simeq H^*(G_T,F).
\]
\end{SGAresult}
\begin{proof}
事实上，由\Term{cohomology}{上同调}的定义，两个\Term{standard-complex}{标准复形}\Term{canonical}{典范}\Term{isomorphism}{同构}。
\end{proof}
\textbf{1.2。\Term{group-morphism}{群态射}的\Term{lifting}{提升}。}\footnote{编者注（39）：本标题为新增。}遵照一般原则，作如下定义。
\begin{SGAresult}{定义}{1.2.1}
设 $1\to M\xrightarrow{u}E\xrightarrow{v}G$ 是一列 $\widehat{\mathcal C}$-\Term{group}{群}的\Term{morphism}{态射}。若下列\Term{equivalence}{等价}条件成立，则称此列\Term{exact-sequence}{正合}：

(i) 对每个 $S\in\operatorname{Ob}\mathcal C$，下列普通\Term{group}{群}的序列\Term{exact-sequence}{正合}：
\[
 1\longrightarrow M(S)\xrightarrow{u(S)}E(S)\xrightarrow{v(S)}G(S).
\]
(ii) 对 $\widehat{\mathcal C}$ 的每个\Term{object}{对象} $H$，下列普通\Term{group}{群}的序列\Term{exact-sequence}{正合}：
\[
 1\longrightarrow\operatorname{Hom}(H,M)\xrightarrow{u(H)}
 \operatorname{Hom}(H,E)\xrightarrow{v(H)}\operatorname{Hom}(H,G).
\]
\end{SGAresult}
在 (ii) 中特别取 $H=G$，可见 $v$ 的\Term{section}{截面}（事先不要求保持\Term{group-structure}{群结构}）构成的\Term{set}{集合}，或者为空，或者是 $\operatorname{Hom}(G,M)$ 下的\Term{principal-homogeneous-set}{主齐性集合}。设其\Term{nonempty}{非空}，并取 $v$ 的一个\Term{section}{截面}
\[
 s:G\longrightarrow E.
\]
则对每个 $S\in\operatorname{Ob}\mathcal C$ 及 $x\in G(S)$，元素 $s(x)\in E(S)$ 定义 $E_S$ 的一个\Term{inner-automorphism}{内自同构}，它使 $M_S$（更准确地说，使 $M_S$ 在 $u_S$ 下的像）保持\Term{invariant}{不变}，因而给出 $M_S$ 的一个\Term{automorphism}{自同构}。
\SourcePage{131}
\begin{SGAresult}{说明}{1.2.1.1}
\footnote{编者注（40）：增加了编号 1.2.1.1 和 1.2.1.2，以突出其中引入的概念。}若 $M$ 交换，从\Term{set-sheaf}{集合层}面可见，该\Term{automorphism}{自同构}不依赖所选\Term{section}{截面}，而只依赖 $x$，且对 $x$ 保持\Term{multiplication}{乘法}。总之，每个 $M$ 交换且 $v$ 有\Term{section}{截面}的\Term{exact-sequence}{正合列}
\[
 (E)\qquad 1\longrightarrow M\xrightarrow{u}E\xrightarrow{v}G
\]
都\Term{adjunction}{伴随}一个 $\widehat{\mathcal C}$-\Term{group}{群}的\Term{morphism}{态射}
\[
 G\longrightarrow\underline{\operatorname{Aut}}_{\widehat{\mathcal C}\text{-gr.}}(M),
\]
称为\Term{extension}{扩张} $(E)$ 定义的 $G$ 对 $M$ 的\Term{action}{作用}。
\end{SGAresult}
\begin{SGAresult}{定义}{1.2.1.2}
第 I 讲 2.3.7 已说明，$v$ 存在一个同时是 $\widehat{\mathcal C}$-\Term{group}{群}的\Term{morphism}{态射}的\Term{section}{截面}，当且仅当\Term{extension}{扩张} $(E)$（“作为\Term{extension}{扩张}”）\Term{isomorphism}{同构}于相对于上述\Term{action}{作用}的 $M$ 与 $G$ 的\Term{semidirect-product}{半直积}。这样的 $v$ 的\Term{section}{截面}称为\Term{extension}{扩张} $(E)$ 的\Term{section}{截面}，或简称 $(E)$ 的\Term{section}{截面}。
\end{SGAresult}
若 $s$ 是 $(E)$ 的一个\Term{section}{截面}，$m\in\Gamma(M)\simeq\operatorname{Ker}(\Gamma(E)\to\Gamma(G))$（$\Gamma$ 的定义见第 I 讲 1.2），则由
\[
 x\longmapsto u(m)s(x)u(m)^{-1}
\]\footnote{编者注（41）：右端的 $x$ 已改为 $s(x)$。}
定义的\Term{morphism}{态射} $G\to E$ 也是 $(E)$ 的\Term{section}{截面}，称其由 $s$ 经 $m$（或 $u(m)$）定义的\Term{inner-automorphism}{内自同构}得到。
\begin{SGAresult}{\Term{lemma}{引理}}{1.2.2}
设 $(E):1\to M\xrightarrow{u}E\xrightarrow{v}G$ 是 $\widehat{\mathcal C}$-\Term{group}{群}的\Term{exact-sequence}{正合列}，其中 $M$ 交换，且 $v$ 有\Term{section}{截面}。使 $G$ 按 $(E)$ 所定义的方式\Term{action}{作用}于 $M$。

(i) \Term{extension}{扩张} $(E)$ \Term{canonical}{典范}地定义一个类 $c(E)\in H^2(G,M)$；此类为零，是 $(E)$ 存在\Term{section}{截面}的充分必要条件。

(ii) 若 $c(E)=0$，则 $(E)$ 的\Term{section}{截面}构成 $Z^1(G,M)$ 下的\Term{principal-homogeneous-set}{主齐性集合}；将\Term{section}{截面}按 $\Gamma(M)$ 中元素定义的\Term{inner-automorphism}{内自同构}的\Term{action}{作用}取商，所得\Term{set}{集合}是 $H^1(G,M)$ 下的\Term{principal-homogeneous-set}{主齐性集合}。

(iii)\footnote{编者注（42）：增加了此项，与 1.2.4 (iii) 对应。另外，此断言对 $v$ 的任意\Term{section}{截面}均成立，见证明。}设 $s$ 为 $(E)$ 的一个\Term{section}{截面}；它在 $\Gamma(M)$ 定义的\Term{inner-automorphism}{内自同构}下的所有\Term{conjugacy}{共轭}，与 $\Gamma(M)/H^0(G,M)$ 之间存在\Term{bijection}{双射}。
\end{SGAresult}
\begin{proof}
证明与普通\Term{group}{群}的情形完全一样；以 $v$ 的一个\Term{section}{截面}为起点，保证了\Term{set}{集合}运算的\Term{functorial}{函子性}。简述其主要步骤如下。

 a) 对 $v$ 的每个\Term{section}{截面} $s$，配以\Term{morphism}{态射}
\[
 Ds:G\times G\longrightarrow M,
\]
\SourcePage{132}
它在\Term{set-sheaf}{集合层}面由下式定义：
\[
 u(Ds(x,y))=s(xy)s(y)^{-1}s(x)^{-1}.
\]
由以下计算可证 $Ds$ 是一个 $2$-\Term{cocycle}{上闭链}。\footnote{编者注（43）：已修改原文，使其与第 I 讲 5.1 相容。}根据\Term{standard-complex}{标准复形}的\Term{differential}{微分}定义（第 I 讲 5.1），有
\[
 (\partial^2Ds)(x,y,z)
 =\bigl(s(x)Ds(y,z)s(x)^{-1}\bigr)\cdot Ds(x,y)^{-1}
 \cdot Ds(xy,z)^{-1}\cdot Ds(x,yz).
\]
只须将 $Ds$ 的定义代入此式，即可得到 $Ds\in Z^2(G,M)$（不需使用任何\Term{commutativity}{交换性}）。

 b) 若 $s,s'$ 是 $v$ 的两个\Term{section}{截面}，则存在 $f:G\to M$，使 $s(x)=f(x)s'(x)$。于是
\begin{align*}
 Ds'(x,y)&=f^{-1}(xy)Ds(x,y)s(x)f(y)s(x)^{-1}f(x),\\
 &=Ds(x,y)\cdot f^{-1}(xy)\cdot\bigl(s(x)f(y)s(x)^{-1}\bigr)\cdot f(x),
\end{align*}
即
\[
 Ds'=Ds\cdot\partial^1f.
\]
\footnote{编者注（44）：下面一句为新增。}这表明 $Ds$ 在 $H^2(G,M)$ 中的类不依赖于 $v$ 的\Term{section}{截面} $s$ 的选取；它就是\Term{extension}{扩张} $(E)$ 的类 $c(E)$。

 c) 设 $s,s'$ 是 $v$ 的两个\Term{section}{截面}，$m\in\Gamma(M)$。等式 $s(x)=m^{-1}s'(x)m$（对每个 $S\in\operatorname{Ob}\mathcal C$ 及 $x\in G(S)$）\Term{equivalence}{等价}于
\[
 s(x)=m^{-1}s'(x)m\,s'(x)^{-1}s'(x),\qquad\text{即}\quad s=\partial^0m\cdot s'.
\]
特别地，$s$ 在 $\Gamma(M)$ 中的\Term{stabilizer}{稳定子}，是满足 $\partial^0m=e_M$ 的 $m\in\Gamma(M)$ 所成的\Term{subgroup}{子群}，也就是\Term{subgroup}{子群} $H^0(G,M)$。这已经证明了 (iii)。

 d) 现在采用通常的论证。\footnote{编者注（45）：下面的论证已详细展开。}取 $v$ 的任意\Term{section}{截面} $s_0$；存在一个为\Term{group-morphism}{群态射}的\Term{section}{截面} $s$，即满足 $Ds=0$ 的\Term{section}{截面}（它必形如 $s=f\cdot s_0$），当且仅当 $(Ds_0)^{-1}=\partial^1f$，也就是当且仅当类 $c(E)$ 为零。这证明了 (i)。

此时，$(E)$ 的\Term{section}{截面}恰为 $s'=h\cdot s$，其中 $h:G\to M$ 满足 $\partial^1h=0$，即 $h\in Z^1(G,M)$。又由 c)，两个这样的\Term{section}{截面} $h_1\cdot s,h_2\cdot s$ 在 $\Gamma(M)$ 下\Term{conjugacy}{共轭}，当且仅当 $h_1,h_2$ 在 $H^1(G,M)$ 中的像相同。这证明了 (ii)。
\end{proof}
仍设
\[
 (E)\qquad 1\longrightarrow M\xrightarrow{u}E\xrightarrow{v}G
\]
是 $\widehat{\mathcal C}$-\Term{group}{群}的\Term{exact-sequence}{正合列}，$M$ 交换。设
\[
 f:H\longrightarrow G
\]
为 $\widehat{\mathcal C}$-\Term{group}{群}的\Term{morphism}{态射}。考虑 $E_f=H\times_G E$；它是一个 $\widehat{\mathcal C}$-\Term{group}{群}，\Term{projection}{投影} $v_f:E_f\to H$ 是 $\widehat{\mathcal C}$-\Term{group}{群}的\Term{morphism}{态射}，$e_f:E_f\to E$ 亦然。另一方面，经 $u$ 将 $M$ 映入 $E$，经\Term{unit-morphism}{单位态射}映入 $H$，便
\SourcePage{133}
定义一个 $\widehat{\mathcal C}$-\Term{group}{群}的\Term{morphism}{态射} $u_f:M\to E_f$。于是构造出 $\widehat{\mathcal C}$-\Term{group}{群}的\Term{commutative-diagram}{交换图}：
\[
\begin{tikzcd}[column sep=large]
 (E)\quad 1\arrow[r]&M\arrow[r,"u"]&E\arrow[r,"v"]&G\\
 (E_f)\quad 1\arrow[r]&M\arrow[u,"\operatorname{id}"]\arrow[r,"u_f"]&E_f\arrow[u,"e_f"]\arrow[r,"v_f"]&H\arrow[u,"f"] .
\end{tikzcd}
\]
立即有：
\begin{SGAresult}{\Term{lemma}{引理}}{1.2.3}
(i) 序列 $(E_f)$ \Term{exact-sequence}{正合}。

(ii) 映射 $s\mapsto e_f\circ s=f'$，在满足 $v_f\circ s=\operatorname{id}$ 的\Term{morphism}{态射} $s:H\to E_f$（即 $v_f$ 的\Term{section}{截面}）与满足 $v\circ f'=f$ 的\Term{morphism}{态射} $f':H\to E$（即“\Term{lifting}{提升}”$f$ 的\Term{morphism}{态射} $f'$）之间给出\Term{bijection}{一一对应}。

(iii) 在上述对应下，$(E_f)$ 的\Term{section}{截面}与\Term{lifting}{提升} $f$ 的\Term{group-morphism}{群态射} $f'$ 互相对应。
\end{SGAresult}
将\Term{lemma}{引理} 1.2.2 应用于\Term{extension}{扩张} $(E_f)$，并结合 1.2.3，得到以下命题（它在\Term{form}{形式}上包含 1.2.2）。
\begin{SGAresult}{命题}{1.2.4}
设 $(E):1\to M\to E\xrightarrow{v}G$ 是 $\widehat{\mathcal C}$-\Term{group}{群}的\Term{exact-sequence}{正合列}，$M$ 交换。设 $f:H\to G$ 是 $\widehat{\mathcal C}$-\Term{group}{群}的\Term{morphism}{态射}，并假设它可\Term{lifting}{提升}为一个（未必为\Term{group-morphism}{群态射}的）\Term{morphism}{态射} $f':H\to E$。使 $H$ 经以下复合\Term{morphism}{态射}\Term{action}{作用}于 $M$（该复合保持\Term{multiplication}{乘法}，且不依赖 $f'$ 的选取）：
\[
 H\xrightarrow{f'}E\xrightarrow{\operatorname{int}}
 \underline{\operatorname{Aut}}_{\widehat{\mathcal C}\text{-gr.}}(M).
\]
(i) \Term{morphism}{态射} $f$ \Term{canonical}{典范}地定义类 $c(f)\in H^2(H,M)$；它为零，是存在\Term{lifting}{提升} $f$ 的 $\widehat{\mathcal C}$-\Term{group}{群}的\Term{morphism}{态射} $f':H\to E$ 的充分必要条件。

(ii) 若 $c(f)=0$，则\Term{lifting}{提升} $f$ 的 $\widehat{\mathcal C}$-\Term{group}{群}的\Term{morphism}{态射} $f'$ 构成的\Term{set}{集合}，按 $\Gamma(M)$ 的元素（即满足 $v(m)=e$ 的 $m\in\Gamma(E)$）定义的\Term{inner-automorphism}{内自同构}的\Term{action}{作用}取商后，是 $H^1(H,M)$ 下的\Term{principal-homogeneous-set}{主齐性集合}。

(iii) 若 $f':H\to E$ 是\Term{lifting}{提升} $f$ 的\Term{group-morphism}{群态射}，则其经 $\Gamma(M)$ 中元素定义的\Term{inner-automorphism}{内自同构}得到的所有变换，构成的\Term{set}{集合}\Term{isomorphism}{同构}于 $\Gamma(M)/\Gamma(M^H)=\Gamma(M)/H^0(H,M)$。
\end{SGAresult}
\SourcePage{134}
\textbf{1.3。\Term{group-law}{群运算}的\Term{group-extension}{扩张}。}考虑如下情形：给定 $\widehat{\mathcal C}$ 的一个\Term{morphism}{态射}
\[
 p:X\longrightarrow Y \tag{$\dagger$}
\]
及一个\Term{action}{作用}于 $X$ 的交换 $\widehat{\mathcal C}$-\Term{group}{群} $M$，使 $X$ 在 $Y$ 上为 $M_Y$ 下的\Term{formally-principal-homogeneous}{形式主齐性对象}。

若 $g:Y\to Z$ 为 $\widehat{\mathcal C}$ 的任意\Term{morphism}{态射}，则 $g\circ p:X\to Z$ 在 $M$ 下\Term{invariant}{不变}：对每个 $S\in\operatorname{Ob}\mathcal C$，$(g\circ p)(S)$ 在 $M(S)$ 对 $X(S)$ 的\Term{action}{作用}下\Term{invariant}{不变}。反过来，对 $n=1,2,3,4$，我们假定以下条件成立。

$(+)_n$：每个在 $M^n$ 对 $X^n$ 的\Term{action}{作用}下\Term{invariant}{不变}的\Term{morphism}{态射} $X^n\to M$，都唯一地通过 $p^n:X^n\to Y^n$ \Term{factorization}{分解}（这里的 $n$ 次幂指\Term{cartesian-power}{笛卡儿幂}）。
\begin{SGAresult}{\Term{lemma}{引理}}{1.3.1}
(i) 若 $h$ 是从 $Y$ 到 $M$ 的\Term{morphism}{态射}，则在\Term{set-sheaf}{集合层}面由 $x\mapsto h(p(x))\cdot x$ 定义的 $X$ 的\Term{automorphism}{自同构} $u_h$ 保持 $p$ 的\Term{fiber}{纤维}，并与 $M$ 对 $X$ 的\Term{action}{作用}交换。\footnote{编者注（46）：这里及下文将“与 $p$ 及 $M$ 的\Term{action}{作用}交换”改为“保持 $p$ 的\Term{fiber}{纤维}，并与 $M$ 的\Term{action}{作用}交换”。此外，增加了以下等式。}也就是说，对所有 $S\in\operatorname{Ob}\mathcal C$、$x\in X(S)$、$m\in M(S)$，有
\[
 p\bigl(h(p(x))\cdot x\bigr)=p(x),\qquad
 m\cdot h(p(x))\cdot x=h(p(m\cdot x))\cdot m\cdot x.
\]
(ii) 这个构造在从 $Y$ 到 $M$ 的\Term{morphism}{态射}与保持 $p$ 的\Term{fiber}{纤维}且与 $M$ 的\Term{action}{作用}交换的 $X$ 的\Term{automorphism}{自同构}之间，建立\Term{bijection}{一一对应}。
\end{SGAresult}
\begin{proof}
第一部分显然，因为 $p(m\cdot x)=p(x)$ 且 $M$ 交换。反过来，保持 $p$ 的\Term{fiber}{纤维}的 $X$ 的\Term{automorphism}{自同构} $u$，在\Term{set-sheaf}{集合层}面可写成 $x\mapsto g(x)x$，其中 $g$ 是某个从 $X$ 到 $M$ 的\Term{morphism}{态射}。若 $u$ 与 $M$ 的\Term{action}{作用}交换，则\Term{morphism}{态射} $g$ 在 $M$ 下\Term{invariant}{不变}\footnote{编者注（47）：事实上，等式 $mg(x)x=mu(x)=u(mx)=g(mx)mx$ 蕴含 $g(mx)=g(x)$。}，再由条件 $(+)_1$ 即得结论。
\end{proof}
现在还假设 $Y$ 上给定一个\Term{group-law}{群运算}，且给定 $Y$ 对 $M$ 的\Term{action}{作用}，也就是 $\widehat{\mathcal C}$-\Term{group}{群}的\Term{morphism}{态射}
\[
 f:Y\longrightarrow\underline{\operatorname{Aut}}_{\widehat{\mathcal C}\text{-gr.}}(M). \tag{$\ddagger$}
\]
\begin{SGAresult}{定义}{1.3.2}
若 $X$ 上的\Term{composition-law}{合成律} $P:X\times X\to X$ 满足以下两个条件，则称其\Term{admissible-law}{可容许}：

(i) $P$ \Term{lifting}{提升} $Y$ 的\Term{group-law}{群运算}，即图
\[
\begin{tikzcd}
 X\times X\arrow[r,"P"]\arrow[d,"{(p,p)}"']&X\arrow[d,"p"]\\
 Y\times Y\arrow[r]&Y
\end{tikzcd}
\]
\SourcePage{135}
交换。

(ii) 对每个 $S\in\operatorname{Ob}\mathcal C$ 以及所有 $x,y\in X(S)$、$m,n\in M(S)$，有关系
\[
 P(m\cdot x,n\cdot y)=m\cdot f(p(x))(n)\cdot P(x,y). \tag{++}
\]
\end{SGAresult}
\begin{SGAresult}{命题}{1.3.3}
$X$ 上的\Term{group-law}{群运算} $*$ \Term{admissible-law}{可容许}，当且仅当下列四个条件成立：

(i) $p:X\to Y$ 为\Term{group-morphism}{群态射}。

(ii) 由 $i(m)=m\cdot e_X$ 定义的\Term{morphism}{态射} $i:M\to X$ 是从 $M$ 到 $\operatorname{Ker}(p)$ 的\Term{group}{群}\Term{isomorphism}{同构}；也就是说，在\Term{set-sheaf}{集合层}面有 $(m\cdot e_X)*(n\cdot e_X)=(mn)\cdot e_X$。

(iii) 对每个 $m\in M(S)$、$x\in X(S)$，有 $m\cdot x=(m\cdot e_X)*x=i(m)*x$。

(iv) $X$ 的\Term{inner-automorphism}{内自同构}按以下\Term{set-sheaf}{集合层}面的公式\Term{action}{作用}于 $\operatorname{Ker}(p)$：
\[
 x*i(m)*x^{-1}=i(f(p(x))m).
\]
\end{SGAresult}
\begin{proof}证明是直接的。\end{proof}
\begin{SGAresult}{\Term{lemma}{引理}}{1.3.4}
\footnote{编者注（48）：为用于 1.3.6 证明的 (e) 项，对陈述及证明作了详细展开。}设 $h:Y\to M$ 为\Term{morphism}{态射}，$u_h$ 为 $X$ 的\Term{automorphism}{自同构} $x\mapsto h(p(x))\cdot x$（见 1.3.1）。设 $P$ 是 $X$ 上\Term{admissible-law}{可容许}的\Term{composition-law}{合成律}（或\Term{group-law}{群运算}），$P'$ 是通过 $u_h$ 从 $P$ 得到的 $X$ 上的\Term{composition-law}{合成律}，即 $P'(x,y)=u_h^{-1}(P(u_h(x),u_h(y)))$。则：

(i) $P'$ 是\Term{admissible-law}{可容许}的\Term{composition-law}{合成律}（或\Term{group-law}{群运算}）。

(ii) 对所有 $x,y\in X(S)$（$S\in\operatorname{Ob}\mathcal C$），令 $v=p(x)$、$w=p(y)$，则
\[
 P'(x,y)=h(vw)^{-1}\cdot h(v)\cdot f(v)(h(w))\cdot P(x,y)
 =(\partial^1h)\bigl(p(x),p(y)\bigr)\cdot P(x,y).
\]
\end{SGAresult}
\begin{proof}
有 $u_h^{-1}=u_{h^\vee}$，其中 $h^\vee:Y\to M$ 由 $h^\vee(y)=h(y)^{-1}$ 定义。根据 1.3.2 (i)、(ii)，有 $P(h(v)\cdot x,h(w)\cdot y)=h(v)\cdot f(v)(h(w))\cdot P(x,y)$ 以及 $p(P(x,y))=vw$，故
\[
 P'(x,y)=h(vw)^{-1}\cdot h(v)\cdot f(v)(h(w))\cdot P(x,y)
 =(\partial^1h)\bigl(p(x),p(y)\bigr)\cdot P(x,y).
\]
由此立即可见 $P'$ 满足 1.3.2 的 (i)、(ii)。
\end{proof}
\begin{SGAresult}{定义}{1.3.5}
若两个\Term{admissible-law}{可容许}\Term{composition-law}{合成律}经由 1.3.4 的方法互相得到，则称它们\Term{equivalence}{等价}。\footnote{编者注（49）：这确实是一个\Term{equivalence-relation}{等价关系}，因为 $u_h^{-1}=u_{h^\vee}$。}
\end{SGAresult}
\begin{SGAresult}{命题}{1.3.6}
假设 $X$ 上存在\Term{admissible-law}{可容许}\Term{composition-law}{合成律}。则：

(i) 存在一个（\Term{canonical}{典范}确定的）类 $c\in H^3(Y,M)$，其为零是 $X$ 上存在结合的\Term{admissible-law}{可容许}\Term{composition-law}{合成律}的充分必要条件。
\SourcePage{136}
(ii) 若 $c=0$，则 $X$ 上结合的\Term{admissible-law}{可容许}\Term{composition-law}{合成律}的\Term{set}{集合}（相应地，这些\Term{composition-law}{合成律}的\Term{equivalence-class}{等价类}的\Term{set}{集合}），是 $Z^2(G,M)$（相应地，$H^2(G,M)$）下的\Term{principal-homogeneous-set}{主齐性集合}。
\end{SGAresult}
\begin{proof}
证明分为若干步骤。

 a) 设 $P$ 是 $X$ 上\Term{admissible-law}{可容许}的\Term{composition-law}{合成律}。因为 $P$ \Term{lifting}{提升} $Y$ 上结合的\Term{composition-law}{合成律}，故存在唯一\Term{morphism}{态射} $a:X^3\to M$，使
\[
 P(x,P(y,z))=a(x,y,z)P(P(x,y),z). \tag{*}
\]
应用条件 1.3.2 (i)、(ii)，立即可见 $a$ 在 $M^3$ 对 $X^3$ 的\Term{action}{作用}下\Term{invariant}{不变}\footnote{编者注（50）：对 $x,y,z\in X(S)$，令 $u=p(x)$、$v=p(y)$、$w=p(z)$。于是
\[
\begin{aligned}
P\bigl(mx,P(m'y,m''z)\bigr)
 &=P\bigl(mx,m'f(v)(m'')P(y,z)\bigr)\\
 &=mf(u)(m')f(uv)(m'')P\bigl(x,P(y,z)\bigr).
\end{aligned}
\]
另一方面，因为 $p(P(x,y))=uv$，也有
\[
\begin{aligned}
P\bigl(P(mx,m'y),m''z\bigr)
 &=P\bigl(mf(u)(m')P(x,y),m''z\bigr)\\
 &=mf(u)(m')f(uv)(m'')P\bigl(P(x,y),z\bigr).
\end{aligned}
\]
将这些等式与 $(*)$ 比较，得到 $a(mx,m'y,m''z)=a(x,y,z)$。}。应用假设 $(+)_3$，得到：

(1) 存在唯一\Term{morphism}{态射} $DP:Y^3\to M$，使
\[
 P(x,P(y,z))=DP(p(x),p(y),p(z))P(P(x,y),z),
\]
而且 $P$ 结合，当且仅当 $DP=0$。

 b) 用前式逐步计算 $P(P(P(x,y),z),t)$。令 $p(x)=u$、$p(y)=v$、$p(z)=w$、$p(t)=h$，得到\footnote{编者注（51）：下图为新增。}下列五边形图；其中\Term{arrow}{箭头} $a\xrightarrow{m}b$ 表示 $b=m\cdot a$：
\[
\begin{tikzcd}[column sep=small,row sep=large]
 &P(x,P(y,P(z,t)))\arrow[dl,"{DP(u,v,wh)}"']\arrow[dr,"{f(u)(DP(v,w,h))}"]&\\
 P(P(x,y),P(z,t))\arrow[d,"{DP(uv,w,h)}"']&&P(x,P(P(y,z),t))\arrow[d,"{DP(u,vw,h)}"]\\
 P(P(P(x,y),z),t)&&P(P(x,P(y,z)),t)\arrow[ll,"{DP(u,v,w)}"]
\end{tikzcd}
\]
因而得到
\[
\begin{aligned}
 {}&DP(u,v,w)\cdot DP(u,vw,h)\cdot f(u)DP(v,w,h)\\
 &\hspace{18mm}\cdot DP(u,v,wh)^{-1}\cdot DP(uv,w,h)^{-1}=e_M,
\end{aligned}
\]
即 $\partial^3DP(u,v,w,h)=e_M$。另一方面，上式左端可用 $P,a$ 写成某个\Term{morphism}{态射} $X^4\to M$ 在 $(x,y,z,t)$ 处的表达式；由 $(+)_4$ 中的唯一性，$\partial^3DP$ 与 $e_M$ 既然\Term{factorization}{分解}同一个\Term{morphism}{态射}，便必相等。因此：

(2) $DP$ 为\Term{cocycle}{上闭链}，即 $DP\in Z^3(Y,M)$。
\SourcePage{137}
 c) 若 $P,P'$ 是 $X$ 上两个\Term{admissible-law}{可容许}\Term{composition-law}{合成律}，则存在唯一\Term{morphism}{态射} $b:X^2\to M$，使 $P'(x,y)=b(x,y)P(x,y)$。对 $P,P'$ 应用 1.3.2 (ii)，可见 $b$ 在 $M^2$ 下\Term{invariant}{不变}，故由 $(+)_2$ 有：

(3) 对每对\Term{admissible-law}{可容许}\Term{composition-law}{合成律} $(P,P')$，存在唯一 $d(P,P'):Y^2\to M$，使
\[
 P'(x,y)=d(P,P')\bigl(p(x),p(y)\bigr)P(x,y),
\]
由此，\Term{admissible-law}{可容许}\Term{composition-law}{合成律}的\Term{set}{集合}成为 $\operatorname{Hom}(Y^2,M)=C^2(Y,M)$ 下的\Term{principal-homogeneous-set}{主齐性集合}。

 d) 在上述条件下，有公式
\[
 DP'-DP=\partial^2d(P,P'). \tag{4}
\]
 e) $P,P'$ \Term{equivalence}{等价}，当且仅当存在\Term{morphism}{态射} $h\in C^1(Y,M)=\operatorname{Hom}(Y,M)$，使 $d(P,P')=\partial^1h$；这由\Term{equivalence}{等价}的定义及 1.3.4 (ii) 得到。

 f) 现在只须归结上述结果。我们要找一个结合的 $P'$，即 $DP'=e_M$。而 $DP$ 是\Term{cocycle}{上闭链}，其在 $H^3(Y,M)$ 中的类不依赖所选\Term{admissible-law}{可容许}\Term{composition-law}{合成律} $P$（由 (3)、(4)）。这个类就是所求的\Term{obstruction}{障碍} $c$。能够选取符合条件的 $P'$，当且仅当 $c=0$；事实上，任取 $P$，由 (1) 需解方程
\[
 0=DP'=DP+\partial^2d(P,P'),
\]
根据 (3)、(4)，这\Term{solvable}{可解}当且仅当 $c=0$。同样由 (3)、(4)，结合的 $P'$ 构成 $Z^2(Y,M)$ 下的\Term{principal-homogeneous-set}{主齐性集合}。按 (e)，这些结合的 $P'$ 在\Term{equivalence}{等价}意义下的\Term{set}{集合}是 $H^2(Y,M)$ 下的\Term{principal-homogeneous-set}{主齐性集合}。
\end{proof}
\SGAsection{2}{群概形态射的无穷小扩张}
沿用第 0 节的记号。设 $Y,X$ 为两个 $S$-\Term{group-functor}{群函子}，$M$ 为\Term{group-morphism}{群态射} $p_X:X\to X^+$ 的\Term{kernel}{核}。于是有 $S$-\Term{group-functor}{群函子}的\Term{exact-sequence}{正合列}
\[
 1\longrightarrow M\longrightarrow X\xrightarrow{p_X}X^+.
\]
由 $X^+$ 的定义，有\Term{isomorphism}{同构}
\[
\begin{aligned}
 \operatorname{Hom}_S(Y,X^+)&\xrightarrow{\sim}\operatorname{Hom}_{S_J}(Y_J,X_J),\\
 \operatorname{Hom}_{S\text{-gr.}}(Y,X^+)&\xrightarrow{\sim}\operatorname{Hom}_{S_J\text{-gr.}}(Y_J,X_J),
\end{aligned}
\]
而\Term{morphism}{态射}
\[
 \operatorname{Hom}_S(Y,p_X):\operatorname{Hom}_S(Y,X)\longrightarrow\operatorname{Hom}_S(Y,X^+)
\]
把 $S$-\Term{morphism}{态射} $f:Y\to X$ 映至 $S$-\Term{morphism}{态射} $f^+:Y\to X^+$；后者由前述\Term{isomorphism}{同构}对应于从 $f$ 经\Term{base-change}{基变换}得到的 $S_J$-\Term{morphism}{态射} $f_J:Y_J\to X_J$。若 $M$ 交换，即可对此情形应用命题 1.2.4。
\SourcePage{138}
\textbf{2.0。}\footnote{编者注（52）：为便于后文引用，增加了编号 2.0。}下文关注以下情形：$Y$ 在 $S$ 上\Term{flat}{平坦}，且 $X$ 是下列两类 $S$-\Term{group-functor}{群函子}之一：

 a) $X$ 是一个 $S$-\Term{group-scheme}{群概形}；

 b) $X=\underline{\operatorname{Aut}}_S(E)$，其中 $E$ 是在 $S$ 上\Term{finite-presentation}{有限表现}的 $S$-\Term{scheme}{概形}。

以 $(\mathrm{Plats})/S$ 表示在 $S$ 上\Term{flat}{平坦}的 $S$-\Term{scheme}{概形}的\Term{category}{范畴}。在情形 (a)（相应地，(b)）中，$S$-\Term{group-functor}{群函子} $M=\operatorname{Ker}(X\to X^+)$、它在 $(\mathrm{Plats})/S$ 上的\Term{restriction}{限制} $L$，以及 $X$ 的\Term{inner-automorphism}{内自同构}对 $M$ 的\Term{action}{作用}，已在 0.9、0.5、0.10（相应地，0.11、0.14、0.12）中算出。也就是说，在情形 (a) 中，设 $L_0$ 是由下式定义的交换 $S_0$-\Term{group-functor}{群函子}：对每个 $S_0$-\Term{scheme}{概形} $T_0$，
\[
 \operatorname{Hom}_{S_0}(T_0,L_0)=
 \operatorname{Hom}_{\mathcal O_{T_0}}
 \bigl(\omega^1_{X_0/S_0}\otimes_{\mathcal O_{S_0}}\mathcal O_{T_0},
 \mathcal J\otimes_{\mathcal O_{S_0}}\mathcal O_{T_0}\bigr),
\]
$X_0$ 经由其在 $\omega^1_{X_0/S_0}$ 中的\Term{adjoint-representation}{伴随表示}\Term{action}{作用}于 $L_0$。则 $L=\prod_{S_0/S}L_0$，即对每个 $S$-\Term{scheme}{概形} $T$，有 $L(T)=L_0(T\times_S S_0)$。

在情形 (b) 中，以 $\pi$ 表示\Term{structure-morphism}{结构态射} $E\to S$；则 $L$ 是 $(\mathrm{Plats})/S$ 上由下式定义的\Term{abelian-group}{交换群}\Term{functor}{函子}：
\[
 \operatorname{Hom}_S(T,L)=
 \Gamma\bigl(T,\pi_*\bigl(\mathscr{H}\!om_{\mathcal O_E}
 (\Omega^1_{E/S},\mathcal J\mathcal O_E)\bigr)\otimes_{\mathcal O_S}\mathcal O_T\bigr),
\]
将 $X$ 视为 $(\mathrm{Plats})/S$ 上的\Term{functor}{函子}时，它按 0.12 所述方式\Term{action}{作用}于 $L$。

于是，在 $(\mathrm{Plats})/S$ 上有\Term{group-functor}{群函子}的\Term{exact-sequence}{正合列}
\[
 (E)\qquad 1\longrightarrow L\longrightarrow X\longrightarrow X^+.
\]
另一方面，因为假设 $Y$ 在 $S$ 上\Term{flat}{平坦}，由 1.1.1，\Term{cohomology}{上同调群} $H^i(Y,M)$ 只依赖 $M$ 在 $(\mathrm{Plats})/S$ 上的\Term{restriction}{限制} $L$。在情形 (a) 中，由于 $L=\prod_{S_0/S}L_0$，故由 1.1.2 有\Term{isomorphism}{同构} $H^i(Y,L)\simeq H^i(Y_0,L_0)$。

综合以上结果，由命题 1.2.4 得到：\footnote{编者注（53）：我们把原来在\Term{theorem}{定理} 2.1 陈述中、构成原书第 113 页主要内容的定义移到了上面。}
\begin{SGAresult}{\Term{theorem}{定理}}{2.1}
设 $S$ 为\Term{scheme}{概形}，$\mathcal I,\mathcal J$ 为两个\Term{quasi-coherent}{拟凝聚}\Term{ideal}{理想}，满足 $\mathcal I\supset\mathcal J$、$\mathcal I\cdot\mathcal J=0$，并定义\Term{closed-subscheme}{闭子概形} $S_0,S_J$。又设：

--- $X$ 为 (a) 或 (b) 型的 $S$-\Term{group-functor}{群函子}，$L_0,L$ 如上；

--- $Y$ 为在 $S$ 上\Term{flat}{平坦}的 $S$-\Term{group-scheme}{群概形}，$f_J:Y_J\to X_J$ 为 $S_J$-\Term{group}{群}的\Term{morphism}{态射}。

则：

(i) $f_J$ 可\Term{lifting}{提升}为 $S$-\Term{group}{群}的\Term{morphism}{态射} $Y\to X$，当且仅当以下两个条件成立：

$(i_1)$ $f_J$ 可\Term{lifting}{提升}为 $S$-\Term{functor}{函子}的\Term{morphism}{态射} $Y\to X$（由 1.2.4，这定义 $Y$ 对 $L$ 的一个\Term{action}{作用}，且不依赖所选\Term{lifting}{提升}；此外，在情形 (a) 中，由此得到的 $Y_0$ 对 $L_0$ 的\Term{action}{作用}来自\Term{morphism}{态射} $f_0:Y_0\to X_0$ 及 $X_0$ 对 $L_0$ 的“\Term{adjoint-action}{伴随作用}”）；
\SourcePage{139}
$(i_2)$ 由 $f_J$ \Term{canonical}{典范}定义的某个\Term{obstruction}{障碍} $c(f_J)$ 为零，其中 $c(f_J)$ 是 $H^2(Y,L)$ 中的一个类（在情形 (a) 中，它\Term{isomorphism}{同构}于 $H^2(Y_0,L_0)$）。

(ii) 若 (i) 的条件成立，则延拓 $f_J$ 的 $S$-\Term{group-functor}{群函子}的\Term{morphism}{态射} $Y\to X$ 构成的\Term{set}{集合} $E$，是 $Z^1(Y,L)$（在情形 (a) 中\Term{isomorphism}{同构}于 $Z^1(Y_0,L_0)$）下的\Term{principal-homogeneous-set}{主齐性集合}；而将 $E$ 按 $X$ 的\Term{inner-automorphism}{内自同构}的\Term{action}{作用}取商，所得\Term{set}{集合}是 $H^1(Y,L)$（在情形 (a) 中\Term{isomorphism}{同构}于 $H^1(Y_0,L_0)$）下的\Term{principal-homogeneous-set}{主齐性集合}，这里所用的\Term{inner-automorphism}{内自同构}由 $X$ 在 $S$ 上、在 $S_J$ 上诱导 $X_J$ 的\Term{unit-section}{单位截面}的\Term{section}{截面}定义。

(iii) 若 $f:Y\to X$ 是延拓 $f_J$ 的 $S$-\Term{group-functor}{群函子}的\Term{morphism}{态射}，则 $f$ 经上述\Term{inner-automorphism}{内自同构}变换所得的所有\Term{morphism}{态射} $Y\to X$，构成的\Term{set}{集合}\Term{isomorphism}{同构}于 $\Gamma(L)/H^0(Y,L)$（在情形 (a) 中\Term{isomorphism}{同构}于 $\Gamma(L_0)/H^0(Y_0,L_0)$）；这里仍取由 $X$ 在 $S$ 上、在 $S_J$ 上诱导 $X_J$ 的\Term{unit-section}{单位截面}的\Term{section}{截面}定义的\Term{inner-automorphism}{内自同构}。
\end{SGAresult}
\textbf{注记 2.1.1。}\footnote{编者注（54）：增加了此注记，作为 4.5.1 的对应，以引入 4.38 使用的记号 $d(f,f')$ 和 $\overline d(f,f')$；因此，也在上面的 2.1 (ii) 中增加了关于 $E$ 本身的部分。}若 $f,f':Y\to X$ 是延拓 $f_J$ 的 $S$-\Term{group-functor}{群函子}的\Term{morphism}{态射}，便得到一个\Term{cocycle}{上闭链} $d(f,f')\in Z^1(Y,L)$（在情形 (a) 中\Term{isomorphism}{同构}于 $Z^1(Y_0,L_0)$），使
\[
 f'=d(f,f')\cdot f. \tag{*}
\]\footnote{编者注（55）：我们沿用了原书的\Term{signature}{符号}约定。本来也可以选写 $f'=d(f',f)\cdot f$，但这样，为在 4.27 中当 $i,i'$ 是两个\Term{immersion}{浸入} $Y\hookrightarrow X$ 时得到等式 $d^1(d(i,i'))=d(Y,i'(Y))$，就须改变 4.5.1 引入的类 $d(Y,Y')$ 的\Term{signature}{符号}，从而导致 4.8、4.14、4.17 公式中的\Term{signature}{符号}变化。我们宁可保留原书的\Term{signature}{符号}（它们全都正确！）。}
以 $\overline d(f,f')$ 表示 $d(f,f')$ 在 $H^1(Y,L)$ 中的像（在情形 (a) 中，此群\Term{isomorphism}{同构}于 $H^1(Y_0,L_0)$）。

\textbf{注记 2.2。}沿用上述记号；特别地，$Y$ 在 $S$ 上\Term{flat}{平坦}。在情形 (b) 中，由 (0.14.1)，$L$ 为\Term{functor}{函子}
\[
 \mathrm W\bigl(\pi_*\bigl(\mathscr{H}\!om_{\mathcal O_E}(\Omega^1_{E/S},\mathcal J\mathcal O_E)\bigr)\bigr)
\]
在 $(\mathrm{Plats})/S$ 上的\Term{restriction}{限制}，其中 $\pi:E\to S$ 为\Term{structure-morphism}{结构态射}。在情形 (a) 中，若还假设 $X$ 在 $S$ 上\Term{locally-finite-presentation}{局部有限表现}，则根据 (0.6.1)，$L$ 为\Term{functor}{函子}
\[
 \prod_{S_0/S}\mathrm W\bigl(\mathscr{H}\!om_{\mathcal O_{S_0}}(\omega^1_{X_0/S_0},\mathcal J)\bigr)
\]
在 $(\mathrm{Plats})/S$ 上的\Term{restriction}{限制}。两种情形中，取 $\mathrm W$ 所用的\Term{module}{模}均为\Term{quasi-coherent}{拟凝聚}，见 EGA I，9.1.1。再假设 $Y$ 在 $S$ 上\Term{affine}{仿射}\footnote{编者注（56）：增加了此假设及第 I 讲 5.3 的引用。}。则根据第 I 讲 5.3，得到
\[
\begin{aligned}
 \text{a)}\quad H^i(Y,L)&=H^i(Y_0,L_0)
 =H^i\bigl(Y_0,\mathscr{H}\!om_{\mathcal O_{S_0}}(\omega^1_{X_0/S_0},\mathcal J)\bigr),\\
 \text{b)}\quad H^i(Y,L)&=H^i\bigl(Y,\pi_*\bigl(\mathscr{H}\!om_{\mathcal O_E}
 (\Omega^1_{E/S},\mathcal J\mathcal O_E)\bigr)\bigr).
\end{aligned}
\]
\SourcePage{140}
\textbf{注记 2.3。}1) 根据 0.16、0.17，当 $Y$ 为\Term{affine}{仿射}\Term{scheme}{概形}且满足下列条件时，条件 $(i_1)$ 自动成立：
\noindent $(*)$ 在情形 (a) 中，$X$ 在 $S$ 上\Term{smooth}{光滑}；在情形 (b) 中，$E$ 在 $S$ 上\Term{smooth}{光滑}且\Term{affine}{仿射}。
2) 此外，在这些条件下（仍假设 $Y$ 在 $S$ 上\Term{flat}{平坦}，见 2.0），在情形 (a) 中，由 2.2 a) 及 (0.6.2)，可写为
\[
 H^i(Y,L)=H^i(Y_0,L_0)
 =H^i\bigl(Y_0,\operatorname{Lie}(X_0/S_0)\otimes_{\mathcal O_{S_0}}\mathcal J\bigr),
\]
\footnote{编者注（57）：以下内容为新增，见 0.14 的编者注（31）。}在情形 (b) 中，由 (0.14.2)、1.1.2 及第 I 讲 5.3，有
\[
 H^i(Y,L)=H^i\bigl(Y_0,\pi_{0*}\mathscr{H}\!om_{\mathcal O_{E_0}}
 (\Omega^1_{E_0/S_0},\mathcal J\otimes_{\mathcal O_{S_0}}\mathcal O_{E_0})\bigr).
\]
现在给出若干推论，讨论 $Y$ 为\Term{diagonalizable-group}{可对角化} $S$-\Term{group}{群}（第 I 讲 4.4）的情形。已知（同讲 5.3.3），若 $S$ \Term{affine}{仿射}，则对 $n>0$ 及任意\Term{quasi-coherent}{拟凝聚} $\mathcal O_S$-\Term{module}{模} $F$，有 $H^n(Y,F)=0$。先看一个特例。
\begin{SGAresult}{推论}{2.4}
设 $S$ 为\Term{scheme}{概形}，$S_0$ 为一个\Term{nilpotent}{幂零}\Term{ideal}{理想}定义的\Term{closed-subscheme}{闭子概形}。设 $Y$ 为\Term{diagonalizable-group}{可对角化} $S$-\Term{group}{群}，且有：

 a) $X$ 为在 $S$ 上\Term{locally-finite-presentation}{局部有限表现}的 $S$-\Term{group}{群}；

 b) $X=\underline{\operatorname{Aut}}_S(E)$，其中 $E$ 为\Term{locally-finite-presentation}{局部有限表现}的 $S$-\Term{scheme}{概形}。

设 $f:Y\to X$ 为 $S$-\Term{group}{群}的\Term{morphism}{态射}，其经\Term{base-change}{基变换}得到的\Term{morphism}{态射} $f_0:Y_0\to X_0$ 是\Term{unit-morphism}{单位态射}。则 $f$ 是\Term{unit-morphism}{单位态射}。
\end{SGAresult}
\begin{proof}
问题在 $S$ 上是\Term{local-scheme}{局部}的，在 (b) 中也在 $E$ 上是\Term{local-scheme}{局部}的。因此可假设 $S$ \Term{affine}{仿射}，且在 (b) 中 $E$ 在 $S$ 上\Term{finite-presentation}{有限表现}。引入由定义 $S_0$ 的\Term{ideal}{理想}的各次幂所定义的 $S$ 的\Term{closed-subscheme}{闭子概形} $S_n$，便归结到 $S_0$ 由\Term{square-zero-ideal}{平方零理想}定义的情形；此时由\Term{theorem}{定理}结合 2.2 即得所述断言。
\end{proof}
不再要求 $f_0$ 必为\Term{unit-morphism}{单位态射}时，有：
\begin{SGAresult}{推论}{2.5}
设 $S,S_0$ 如 2.4，再假设 $S$ \Term{affine}{仿射}。设 $Y$ 为\Term{diagonalizable-group}{可对角化} $S$-\Term{group}{群}，$X$ 为 $S$-\Term{group-functor}{群函子}，$f_0:Y_0\to X_0$ 为 $S_0$-\Term{group-functor}{群函子}的\Term{morphism}{态射}。

(i)\footnote{编者注（58）：原书的陈述为：“假设下列两项之一成立：(a) $X$ 是\Term{locally-finite-presentation}{局部有限表现}的 $S$-\Term{group}{群}；(b) $X=\underline{\operatorname{Aut}}_S(E)$，其中 $E$ 在 $S$ 上\Term{finite-presentation}{有限表现}。则 $f_0$ 可延拓为 $S$-\Term{group}{群}的\Term{morphism}{态射} $Y\to X$，当且仅当它可延拓为 $S$-\Term{functor}{函子}的\Term{morphism}{态射} $Y\to X$。”并说明：“(i) 由\Term{theorem}{定理}的 (ii) 部分逐步\Term{pushout}{推出}。”这个证明似乎不够充分：例如若 $\mathcal I^3=0$、$\mathcal J=\mathcal I^2$，且 $f:Y\to X$ 为\Term{lifting}{提升} $f_0$ 的 $S$-\Term{functor}{函子}的\Term{morphism}{态射}，则 $f_0$ 可\Term{lifting}{提升}为一个 $S_J$-\Term{group}{群}的\Term{morphism}{态射} $g_J:Y_J\to X_J$；但随后，$g_J$ 是否又可\Term{lifting}{提升}为 $S$-\Term{functor}{函子}的\Term{morphism}{态射} $g:Y\to X$ 呢？无论如何，下文未使用断言 (i)，其中始终假设 $X$ 在 $S$ 上\Term{smooth}{光滑}。}
\SourcePage{141}
(ii) 假设以下两项之一成立：

(a) $X$ 为在 $S$ 上\Term{smooth}{光滑}的 $S$-\Term{group}{群}；

(b) $X=\underline{\operatorname{Aut}}_S(E)$，其中 $E$ 在 $S$ 上\Term{smooth}{光滑}且\Term{affine}{仿射}。

则 $f_0$ 可延拓为 $S$-\Term{group}{群}的\Term{morphism}{态射} $Y\to X$；任意两个这样的延拓，都经由 $X$ 的一个\Term{inner-automorphism}{内自同构}\Term{conjugacy}{共轭}，该\Term{automorphism}{自同构}由 $X$ 在 $S$ 上、在 $S_0$ 上诱导 $X_0$ 的\Term{unit-section}{单位截面}的\Term{section}{截面}定义。
\end{SGAresult}
\begin{proof}
如上引入各 $S_n$。\footnote{编者注（59）：以下内容对原书作了少许修改并详细展开。}对于 (ii)，先注意在 $S$ 上\Term{smooth}{光滑}的\Term{scheme}{概形}必在 $S$ 上\Term{locally-finite-presentation}{局部有限表现}；故在情形 (b) 中，$E$ 在 $S$ 上\Term{smooth}{光滑}且\Term{affine}{仿射}，必在 $S$ 上\Term{finite-presentation}{有限表现}，即确实满足 2.0 的假设 (b)。

于是，在 (ii) 的假设下，由 0.16、0.17，2.1 的条件 $(i_1)$ 自动成立；此外，$X_{S_n}$ 在 $S_n$ 上的任意\Term{section}{截面}，均可\Term{lifting}{提升}为 $X_{S_{n+1}}$ 在 $S_{n+1}$ 上的\Term{section}{截面}：在情形 (a) 中由“在 $S$ 上\Term{smooth}{光滑}”的定义，在情形 (b) 中由 0.17。因此，若 $f,f'$ 是 $f_0$ 的两个\Term{lifting}{提升}，则逐次\Term{lifting}{提升}\Term{theorem}{定理} (ii) 所保证存在的\Term{inner-automorphism}{内自同构}，可逐步假设 $f_n=f'_n$。证明完成。
\end{proof}
同样论证，并结合注记 2.3，得到：
\begin{SGAresult}{推论}{2.6}
设 $S$ 为\Term{scheme}{概形}，$\mathcal I$ 为定义\Term{closed-subscheme}{闭子概形} $S_0$ 的\Term{nilpotent}{幂零}\Term{ideal}{理想}；$Y$ 为在 $S$ 上\Term{flat}{平坦}且\Term{affine}{仿射}的 $S$-\Term{group}{群}，$X$ 为在 $S$ 上\Term{smooth}{光滑}的 $S$-\Term{group}{群}。

(i) 若对每个 $n\geq0$ 均有 $H^2(Y_0,\operatorname{Lie}(X_0/S_0)\otimes_{\mathcal O_{S_0}}\mathcal I^{n+1}/\mathcal I^{n+2})=0$，则每个 $S_0$-\Term{group}{群}的\Term{morphism}{态射} $f_0:Y_0\to X_0$ 均可\Term{lifting}{提升}为 $S$-\Term{group}{群}的\Term{morphism}{态射} $f:Y\to X$。

(ii) 若对每个 $n\geq0$ 均有 $H^1(Y_0,\operatorname{Lie}(X_0/S_0)\otimes_{\mathcal O_{S_0}}\mathcal I^{n+1}/\mathcal I^{n+2})=0$，则两个这样的\Term{lifting}{提升}经由 $X$ 的一个\Term{inner-automorphism}{内自同构}\Term{conjugacy}{共轭}，该\Term{automorphism}{自同构}由 $X$ 在 $S$ 上、在 $S_0$ 上诱导 $X_0$ 的\Term{unit-section}{单位截面}的\Term{section}{截面}定义。
\end{SGAresult}
有如下\Term{lemma}{引理}。
\begin{SGAresult}{\Term{lemma}{引理}}{2.7}
设 $S$ 为\Term{affine}{仿射}\Term{scheme}{概形}，$G$ 为\Term{affine}{仿射} $S$-\Term{group}{群}，$\mathcal F$ 为\Term{quasi-coherent}{拟凝聚} $\mathcal O_S$-\Term{module}{模}，$\mathcal L$ 为\Term{locally-free}{局部自由} $\mathcal O_S$-\Term{module}{模}。假设按第 I 讲的意义给定 $G$ 对 $\mathcal F$ 的\Term{action}{作用}，这定义 $G$ 对 $\mathcal F\otimes_{\mathcal O_S}\mathcal L$ 的\Term{action}{作用}\footnote{编者注（60）：这里令 $\mathcal L$ 带 $G$ 的\Term{trivial-action}{平凡作用}。}。以 $\Lambda$ 表示 $S$ 的\Term{ring}{环}，$L$ 表示定义 $\mathcal L$ 的 $\Lambda$-\Term{module}{模}（因而它是\Term{projective-module}{投射模}）。有\Term{canonical}{典范}\Term{isomorphism}{同构}
\[
 H^*(G,\mathcal F\otimes_{\mathcal O_S}\mathcal L)
 \simeq H^*(G,\mathcal F)\otimes_\Lambda L.
\]
\end{SGAresult}
\begin{proof}
\footnote{编者注（61）：我们详细展开并简化了原书的证明（原书还调用了\Term{isomorphism}{同构} $\mathcal H^n(\mathcal C\otimes_{\mathcal O_S}\mathcal L)\simeq\mathcal H^n(\mathcal C)\otimes_{\mathcal O_S}\mathcal L$ 及 $H^n(\Gamma(-))\simeq\Gamma(\mathcal H^n(-))$，其中 $\mathcal H^n(\mathcal C)$ 表示\Term{complex}{复形} $\mathcal C$ 的\Term{cohomology-sheaf}{上同调层}）。}以 $\mathcal A$ 表示 $\mathcal O_S$-\Term{algebra}{代数} $\mathcal A(G)$，考虑\Term{quasi-coherent}{拟凝聚} $\mathcal O_S$-\Term{module}{模}的\Term{complex}{复形} $\mathcal C$：
\[
 0\longrightarrow\mathcal F\longrightarrow\mathcal F\otimes_{\mathcal O_S}\mathcal A
 \longrightarrow\mathcal F\otimes_{\mathcal O_S}\mathcal A\otimes_{\mathcal O_S}\mathcal A\longrightarrow\cdots
\]
\SourcePage{142}
根据第 I 讲 5.3，$H^*(G,\mathcal F)$（相应地，$H^*(G,\mathcal F\otimes_{\mathcal O_S}\mathcal L)$）是\Term{complex}{复形} $\Gamma(S,\mathcal C)$（相应地，$\Gamma(S,\mathcal C\otimes_{\mathcal O_S}\mathcal L)$）的\Term{cohomology}{上同调}。因为 $S$ \Term{affine}{仿射}，有（见 EGA I，1.3.12）
\[
 \Gamma(S,\mathcal C\otimes_{\mathcal O_S}\mathcal L)
 \simeq\Gamma(S,\mathcal C)\otimes_\Lambda L.
\]
又因 $L$ 是\Term{projective-module}{投射}（因而\Term{flat}{平坦}）$\Lambda$-\Term{module}{模}，亦有
\[
 H^*\bigl(\Gamma(S,\mathcal C)\otimes_\Lambda L\bigr)
 \simeq H^*(\Gamma(S,\mathcal C))\otimes_\Lambda L,
\]
所述结果随之得到。
\end{proof}
由此\Term{lemma}{引理}，2.6 转化为：
\begin{SGAresult}{推论}{2.8}
设 $S$ 为\Term{affine}{仿射}\Term{scheme}{概形}，$\mathcal I$ 为 $S$ 上定义\Term{closed-subscheme}{闭子概形} $S_0$ 的\Term{nilpotent}{幂零}\Term{ideal}{理想}。假设 $\mathcal I^{n+1}/\mathcal I^{n+2}$ 在 $S_0$ 上\Term{locally-free}{局部自由}。设 $Y$ 为在 $S$ 上\Term{flat}{平坦}且\Term{affine}{仿射}的 $S$-\Term{group}{群}，$X$ 为在 $S$ 上\Term{smooth}{光滑}的 $S$-\Term{group}{群}，$f_0:Y_0\to X_0$ 为 $S$-\Term{group}{群}的\Term{morphism}{态射}。

(i) 若 $H^2(Y_0,\operatorname{Lie}(X_0/S_0))=0$，则 $f_0$ 可\Term{lifting}{提升}为 $S$-\Term{group}{群}的\Term{morphism}{态射} $Y\to X$。

(ii) 若 $H^1(Y_0,\operatorname{Lie}(X_0/S_0))=0$，则两个这样的\Term{lifting}{提升}经由 $X$ 的一个\Term{inner-automorphism}{内自同构}\Term{conjugacy}{共轭}，该\Term{automorphism}{自同构}由 $X$ 在 $S$ 上、在 $S_0$ 上诱导 $X_0$ 的\Term{unit-section}{单位截面}的\Term{section}{截面}定义。
\end{SGAresult}
特别地，取 $Y=X$，有：
\begin{SGAresult}{推论}{2.9}
设 $S,S_0$ 如上，$X$ 为在 $S$ 上\Term{smooth}{光滑}且\Term{affine}{仿射}的 $S$-\Term{group}{群}。

(i) 若 $H^1(X_0,\operatorname{Lie}(X_0/S_0))=0$，则 $X$ 在 $S$ 上的每个在 $X_0$ 上诱导\Term{identity-morphism}{恒等态射}的\Term{endomorphism}{自同态}，都是由 $X$ 在 $S$ 上、在 $S_0$ 上诱导 $X_0$ 的\Term{unit-section}{单位截面}的\Term{section}{截面}定义的\Term{inner-automorphism}{内自同构}。

(ii) 若 $H^2(X_0,\operatorname{Lie}(X_0/S_0))=0$，则 $X_0$ 的任意 $S_0$-\Term{automorphism}{自同构}均可延拓为 $X$ 的 $S$-\Term{automorphism}{自同构}。\footnote{编者注（62）：事实上，设 $f_0$ 为 $X_0$ 的 $S_0$-\Term{automorphism}{自同构}，$g_0$ 为其逆。由 2.8 (i)，$f_0$（相应地，$g_0$）可\Term{lifting}{提升}为 $X$ 的 $S$-\Term{endomorphism}{自同态} $f$（相应地，$g$）。于是 $g\circ f$ 与 $f\circ g$ 均为在 $X_0$ 上诱导\Term{identity-morphism}{恒等态射}的 $X$ 的\Term{endomorphism}{自同态}；故由 0.7，它们都是 $X$ 的 $S$-\Term{automorphism}{自同构}，因而 $f,g$ 亦然。}
\end{SGAresult}
\textbf{注记 2.10。}由\Term{theorem}{定理}，关于 $H^1$ 的断言都有逆命题。举例如下：若 $S=\mathrm I_{S_0}$ 是 $S_0$ 上的\Term{dual-numbers-scheme}{对偶数概形}（第 II 讲 2.1），$X$ 是\Term{flat}{平坦} $S$-\Term{group}{群}，且 $X$ 在 $S$ 上的每个在 $S_0$ 上诱导\Term{identity-morphism}{恒等态射}的\Term{automorphism}{自同构}，都是由 $X$ 在 $S$ 上、在 $S_0$ 上诱导 $X_0$ 的\Term{unit-section}{单位截面}的\Term{section}{截面}定义的\Term{inner-automorphism}{内自同构}，则 $H^1(X_0,\operatorname{Lie}(X_0/S_0))=0$。\footnote{编者注（63）：第 XXIV 讲 1.13 使用了这一点。}
\begin{SGAresult}{推论}{2.11}
设 $S,\mathcal I,\mathcal J$ 如 2.1。设 $Y$ 为在 $S$ 上\Term{flat}{平坦}的 $S$-\Term{group-scheme}{群概形}，$X$ 为 $S$-\Term{group-scheme}{群概形}，$f:Y\to X$ 为 $S$-\Term{group}{群}的\Term{morphism}{态射}。由 $f$ 经那些在 $X(S_J)$ 中诱导\Term{identity-element}{单位元}的 $x\in X(S)$ \Term{conjugacy}{共轭}得到的从 $Y$ 到 $X$ 的\Term{morphism}{态射}，构成的\Term{set}{集合}\Term{isomorphism}{同构}于商
\[
 E=\operatorname{Hom}_{\mathcal O_{S_0}}(\omega^1_{X_0/S_0},\mathcal J)
 \big/\operatorname{Hom}_{\mathcal O_{S_0}}(\omega^1_{X_0/S_0},\mathcal J)^{\operatorname{ad}(Y_0)},
\]
\SourcePage{143}
其中第二个\Term{group}{群}由满足下列条件的 $\mathcal O_{S_0}$-\Term{morphism}{态射} $\omega^1_{X_0/S_0}\to\mathcal J$ 组成：任意\Term{base-change}{基变换} $S'\to S_0$ 后，所得\Term{morphism}{态射}
\[
 \omega^1_{X_0/S_0}\otimes_{\mathcal O_{S_0}}\mathcal O_{S'}
 \longrightarrow\mathcal J\otimes_{\mathcal O_{S_0}}\mathcal O_{S'}
\]
在 $Y_0(S')$ 对第一个因子的\Term{action}{作用}下\Term{invariant}{不变}。
\end{SGAresult}
\begin{proof}
由 2.1 (iii)，所求\Term{set}{集合}\Term{isomorphism}{同构}于 $\Gamma(L_0)/H^0(Y_0,L_0)$。而 $\Gamma(L_0)=\operatorname{Hom}_{\mathcal O_{S_0}}(\omega^1_{X_0/S_0},\mathcal J)$，且 $H^0(Y_0,L_0)$ 显然就是前述意义下的 $\Gamma(L_0)^{\operatorname{ad}(Y_0)}$。
\end{proof}
\begin{SGAresult}{推论}{2.12}
在 2.11 的条件下，若再假设 $\omega^1_{X_0/S_0}$ 是\Term{finite-rank}{有限秩}\Term{locally-free}{局部自由}的，则
\[
 E\simeq\Gamma\bigl(S_0,\operatorname{Lie}(X_0/S_0)\otimes_{\mathcal O_{S_0}}\mathcal J\bigr)
 \big/H^0\bigl(Y_0,\operatorname{Lie}(X_0/S_0)\otimes_{\mathcal O_{S_0}}\mathcal J\bigr).
\]
\end{SGAresult}
\begin{proof}
\footnote{编者注（64）：以下一句为新增。}事实上，若 $\omega^1_{X_0/S_0}$ 是\Term{finite-rank}{有限秩}\Term{locally-free}{局部自由}的，则有
\[
 \mathscr{H}\!om_{\mathcal O_{S_0}}(\omega^1_{X_0/S_0},\mathcal J)
 \simeq\operatorname{Lie}(X_0/S_0)\otimes_{\mathcal O_{S_0}}\mathcal J.
\]
\end{proof}
\begin{SGAresult}{推论}{2.13}
若再假设 $Y_0$ \Term{diagonalizable-group}{可对角化}，则
\[
 E\simeq\Gamma\bigl(S_0,\operatorname{Lie}(X_0/S_0)\otimes_{\mathcal O_{S_0}}\mathcal J\bigr)
 \big/\Gamma\bigl(S_0,\operatorname{Lie}(X_0/S_0)^{\operatorname{ad}(Y_0)}\otimes_{\mathcal O_{S_0}}\mathcal J\bigr),
\]
其中 $\operatorname{Lie}(X_0/S_0)^{\operatorname{ad}(Y_0)}$ \Term{constructible}{可构造}为第 I 讲 4.7.3 的\Term{decomposition}{分解}中，对应 $Y_0$ 的零\Term{character}{特征标}的因子。
\end{SGAresult}
\begin{proof}
事实上，若 $Y_0\simeq\mathrm D_{S_0}(M)$，则由该处有\Term{direct-sum}{直和}\Term{decomposition}{分解}
\[
 \operatorname{Lie}(X_0/S_0)=\operatorname{Lie}(X_0/S_0)_0
 \oplus\bigoplus_{\substack{m\in M\\m\ne0}}\operatorname{Lie}(X_0/S_0)_m.
\]
与 $\mathcal J$ 作\Term{tensor-product}{张量积}，得到 $\operatorname{Lie}(X_0/S_0)\otimes_{\mathcal O_{S_0}}\mathcal J$ 的类似\Term{decomposition}{分解}，故有关系
\[
 H^0\bigl(Y_0,\operatorname{Lie}(X_0/S_0)\otimes\mathcal J\bigr)
 \simeq\Gamma\bigl(S_0,\operatorname{Lie}(X_0/S_0)_0\otimes_{\mathcal O_{S_0}}\mathcal J\bigr).
\]
\end{proof}
\begin{SGAresult}{推论}{2.14}
若再假设 $S_0$ \Term{affine}{仿射}，则
\[
 E\simeq\Gamma\left(S_0,
 \left[\operatorname{Lie}(X_0/S_0)/\operatorname{Lie}(X_0/S_0)^{\operatorname{ad}(Y_0)}\right]
 \otimes_{\mathcal O_{S_0}}\mathcal J\right).
\]
\end{SGAresult}
\SGAsection{3}{群概形的无穷小扩张}
仍沿用第 0 节的记号（$S,\mathcal I,\mathcal J$ 等）。给定一个 $S$-\Term{scheme}{概形} $X$，并假设 $X_J$ 带有\Term{group-structure}{群结构}。我们要找出 $X$ 上在 $X_J$ 上诱导给定结构的 $S$-\Term{group-structure}{群结构}。

从现在起，假设 $X$ 在 $S$ 上\Term{flat}{平坦}。设 $\mathcal C$ 为在 $S$ 上\Term{flat}{平坦}的 $S$-\Term{scheme}{概形}的\Term{category}{范畴}，于是 $X\in\operatorname{Ob}\mathcal C$。分别以 $Y,M$ 表示由 $X^+,L'_X$ 定义的 $\mathcal C$ 上的\Term{functor}{函子}。\Term{canonical}{典范}\Term{morphism}{态射} $p_X:X\to X^+$ 定义 $\widehat{\mathcal C}$ 中的一个\Term{morphism}{态射}
\[
 p:X\longrightarrow Y,
\]
\SourcePage{144}
而 $\widehat{(\mathrm{Sch})}_{/S}$ 中 $L'_X$ 对 $X$ 的\Term{action}{作用}，定义了 $\widehat{\mathcal C}$ 中 $M$ 对 $X$ 的\Term{action}{作用}。立即可见，$X$ 如此确实成为在 $Y$ 上的 $M_Y$ 下的\Term{formally-principal-homogeneous}{形式主齐性对象}（见 0.2 (i) 及 0.4）。

0.8 定义的 $X^+$ 对 $L'_X$ 的\Term{action}{作用}（该处记为 $\operatorname{Ad}$），定义 $Y$ 对 $M$ 的一个\Term{action}{作用}，记为 $f$。另一方面，由 0.5 知
\[
 \operatorname{Hom}_{\widehat{\mathcal C}}(Z,M)
 \simeq\operatorname{Hom}_{S_0}(Z_0,L_0),\qquad Z\in\operatorname{Ob}\mathcal C,
\]
其中 $L_0$ 为 0.5 定义的\Term{functor}{函子}。
\begin{SGAresult}{\Term{lemma}{引理}}{3.1}
(i) 对每个正整数 $n$，1.3 的条件 $(+)_n$ 均成立。

(ii) 若令 $S_0$-\Term{group}{群} $X_0$ 经其\Term{adjoint-representation}{伴随表示}\Term{action}{作用}于 $S_0$-\Term{functor}{函子} $L_0$，则有\Term{canonical}{典范}\Term{isomorphism}{同构}
\[
 H^*(X_0,L_0)\simeq H^*(Y,M)
\]
（第一个\Term{cohomology}{上同调}在 $(\mathrm{Sch})/S_0$ 中计算，第二个在 $\mathcal C$ 中计算）。
\end{SGAresult}
\begin{proof}
\Term{lemma}{引理}的两部分均由以下关系得到：
\[
\begin{aligned}
 \operatorname{Hom}_{\widehat{\mathcal C}}(Y,M)
 &\simeq\operatorname{Hom}_{\widehat{(\mathrm{Sch})}_{/S_0}}(X^+\times_S S_0,L_0)\\
 &\simeq\operatorname{Hom}_{S_0}(X_0,L_0)\\
 &\simeq\operatorname{Hom}_{\widehat{\mathcal C}}(X,M),
\end{aligned}
\]
它直接来自将 $M$ 定义为一个 $\prod_{S_0/S}$。更一般地，用 $X^n,Y^n$ 代替 $X,Y$ 时，此关系仍成立；因此每个\Term{morphism}{态射} $X^n\to M$ 都唯一地通过 $Y^n$ \Term{factorization}{分解}，这蕴含 $(+)_n$。也由此得到关系 $C^*(Y,M)=C^*(X_0,L_0)$，从而得到 (ii)。
\end{proof}
因此可应用 1.3 的各项构造。特别地：
\begin{SGAresult}{\Term{lemma}{引理}}{3.2}
设 $P:X\times_S X\to X$ 为\Term{morphism}{态射}。$P$ 诱导 $X_J$ 的\Term{group-law}{群运算}，当且仅当 $P$ 是 $X$ 上的\Term{admissible-law}{可容许}\Term{composition-law}{合成律}（见 1.3.2）。
\end{SGAresult}
\begin{proof}
事实上，$P$ 诱导 $X_J$ 的\Term{group-law}{群运算}，当且仅当 $P$ \Term{lifting}{提升} $X^+$ 的\Term{group-law}{群运算}，也即 $Y$ 的\Term{group-law}{群运算}。因此只须证明，每个\Term{lifting}{提升} $X_J$ 的\Term{group-law}{群运算}的\Term{morphism}{态射} $P$ 均满足 1.3.2 (ii) 的恒等式 $(++)$；而这恰是 0.8 已证明的内容。
\end{proof}
\begin{SGAresult}{命题}{3.3}
设 $S$ 为\Term{scheme}{概形}，$S_0$ 为由一个\Term{nilpotent}{幂零}\Term{ideal}{理想}定义的\Term{closed-subscheme}{闭子概形}。设 $X$ 为在 $S$ 上\Term{flat}{平坦}且\Term{quasi-compact}{拟紧}或\Term{locally-finite-presentation}{局部有限表现}的 $S$-\Term{scheme}{概形}。设 $P:X\times_S X\to X$ 为 $X$ 上的\Term{composition-law}{合成律}。$P$ 是\Term{group-law}{群运算}，当且仅当下列两个条件成立：

(i) $P$ 结合。

(ii) $P$ 在 $X_0=X\times_X S_0$ 上诱导\Term{group-law}{群运算}。
\end{SGAresult}
\SourcePage{145}
\begin{proof}
这些条件显然必要。以下证明它们充分。先假设 $X\to S$ 有\Term{section}{截面}。此时对每个 $S'\to S$，$X(S')$ 均\Term{nonempty}{非空}，故只须\footnote{编者注（65）：我们修改了原书，删去了对 Bourbaki 关于\Term{semigroup}{半群}的习题的不恰当引用（见 [BAlg]，\S I.2，习题 9 至 13），并说明了左、右\Term{translation}{平移}的\Term{action}{作用}，见下一条编者注。}证明，对每个 $x\in X(S')$，由 $x$ 定义的左、右\Term{translation}{平移}均为 $X_{S'}$ 的\Term{isomorphism}{同构}。\footnote{编者注（66）：设\Term{nonempty}{非空}\Term{set}{集合} $E$ 带有结合的\Term{composition-law}{合成律}，使每个左\Term{translation}{平移} $\ell_x$ 均为\Term{bijection}{双射}；固定 $x_0\in E$。存在唯一 $e\in E$，使 $x_0\cdot e=x_0$；于是由 $x_0\cdot e\cdot x=x_0\cdot x$ 得 $e\cdot x=x$，对所有 $x\in E$ 成立。另一方面，对每个 $x$，存在唯一 $x'$，使 $x\cdot x'=e$。再假设存在 $b\in E$，使右\Term{translation}{平移} $r_b$ 为\Term{injectivity}{单射}。于是对每个 $x$，由等式 $x\cdot e\cdot b=x\cdot b$ 得 $x\cdot e=x$（即 $e$ 是\Term{identity-element}{单位元}）；又由 $x\cdot x'\cdot x=x=x\cdot e$ 得 $x'\cdot x=e$，即 $x'$ 同时为 $x$ 的左逆与右逆，故 $E$ 为\Term{group}{群}。注意，“$r_b$ \Term{injectivity}{单射}”的假设是必要的：在任意\Term{set}{集合} $E$ 上，都可对所有 $x,y\in E$ 定义\Term{composition-law}{合成律} $x\cdot y=y$；此时每个左\Term{translation}{平移}都是\Term{identity-morphism}{恒等映射}（这也给出该\Term{composition-law}{合成律}的\Term{associativity}{结合性}），但对每个 $y$，有 $r_y(E)=\{y\}$，故当 $|E|>1$ 时，$E$ 不是\Term{group}{群}。}

显然可假设 $S'=S$；所讨论的\Term{translation}{平移} $t$ 在 $X_0$ 上诱导 $X_0$ 的一个\Term{translation}{平移} $t_0$；因为 $X_0$ 为\Term{group}{群}，故 $t_0$ 是\Term{automorphism}{自同构}。由\Term{flat}{平坦性}即可得结论（SGA 1，第 III 讲 4.2）。\footnote{编者注（67）：由于 $X$ 与 $X_0$ 有相同的底\Term{topological-space}{拓扑空间}，且 $t_0$ 为\Term{automorphism}{自同构}，故 $t$ 是\Term{homeomorphism}{同胚}，因而是\Term{affine-morphism}{仿射态射}，见第 VI B 讲 2.9.1 或 EGA IV$_4$，18.12.7.1。因此只须证明：若 $J$ 是\Term{ring}{环} $\Lambda$ 的一个\Term{nilpotent}{幂零}\Term{ideal}{理想}，$\varphi:A\to B$ 是 $\Lambda$-\Term{algebra}{代数}的\Term{morphism}{态射}，$B$ 在 $\Lambda$ 上\Term{flat}{平坦}，且 $\varphi\otimes_\Lambda(\Lambda/J)$ 为\Term{bijection}{双射}，则 $\varphi$ 为\Term{bijection}{双射}。由“\Term{nilpotent-nakayama-lemma}{幂零 Nakayama 引理}”，$\varphi$ 为\Term{surjection}{满射}；此外，$B$ 在 $\Lambda$ 上\Term{flat}{平坦}，故亦有 $\operatorname{Ker}(\varphi)\otimes_\Lambda(\Lambda/J)=0$，从而 $\operatorname{Ker}(\varphi)=0$，故 $\varphi$ 为\Term{bijection}{双射}。}

现在不再假设 $X$ 在 $S$ 上有\Term{section}{截面}，而是假设存在 $S'\to S$，使 $X_{S'}$ 在 $S'$ 上有\Term{section}{截面}。由刚才的结论，$X_{S'}$ 是 $S'$-\Term{group}{群}；考虑其\Term{unit-section}{单位截面} $e'$。$e'$ 在 $\operatorname{pr}_i:S'\times_S S'\to S'$（$i=1,2$）下的\Term{inverse-image}{逆像}，是 $X_{S''}$ 对于 $P_{S'}$ 在 $\operatorname{pr}_i$ 下的\Term{inverse-image}{逆像}\Term{group-law}{群运算}的\Term{unit-section}{单位截面}。但因为 $P$ “定义在 $S$ 上”，这两个\Term{group-law}{群运算}相同，故其\Term{unit-section}{单位截面}也相同。因此 $\operatorname{pr}_1^*(e')=\operatorname{pr}_2^*(e')$。

若 $S'\to S$ 是\Term{descent-morphism}{下降态射}（见第 IV 讲第 2 节），便存在 $X$ 的一个\Term{section}{截面}，其经\Term{base-extension}{基扩张}得到 $e'$，从而证明完成。由于 $X_X$ 在 $X$ 上有\Term{section}{截面}（\Term{diagonal-section}{对角截面}），可见现在只须证明 $X\to S$ 是\Term{descent-morphism}{下降态射}。但它\Term{flat}{平坦}且\Term{surjection}{满射}，又\Term{quasi-compact}{拟紧}或\Term{locally-finite-presentation}{局部有限表现}，故为 (fpqc) 的\Term{covering}{覆盖}，因而是\Term{descent-morphism}{下降态射}（第 IV 讲第 6 节）。
\end{proof}
\textbf{注记。}事实上，$X\to S$ \Term{quasi-compact}{拟紧}或\Term{locally-finite-presentation}{局部有限表现}的假设是多余的。这由以下结果得到，读者可将其作为第 IV 讲的习题来证明：

在正文关于 $S,S_0$ 的条件下，若 $X\to S$ 是\Term{flat}{平坦}\Term{morphism}{态射}，且 $X_0\to S_0$ 是 (fpqc) 的\Term{covering}{覆盖}\Term{morphism}{态射}，则 $X\to S$ 是\Term{descent-morphism}{下降态射}。
\SourcePage{146}
\begin{SGAresult}{\Term{lemma}{引理}}{3.4}
$X$ 上两个\Term{admissible-law}{可容许}\Term{composition-law}{合成律}\Term{equivalence}{等价}（见 1.3.5），当且仅当它们经由 $X$ 在 $S$ 上的一个在 $X_J$ 上诱导\Term{identity-morphism}{恒等态射}的\Term{automorphism}{自同构}互相得到。
\end{SGAresult}
\begin{proof}
事实上，1.3.1 中构造的\Term{morphism}{态射}，恰好就是上述陈述中的那些\Term{morphism}{态射}（由 0.7）。\footnote{编者注（68）：事实上，根据 0.7 的证明，$X$ 的那些在 $X_J$ 上诱导\Term{identity-morphism}{恒等态射}的 $S$-\Term{endomorphism}{自同态}，就是\Term{automorphism}{自同构} $m\cdot\operatorname{id}_X$，其中 $m$ 遍历 $M(X)=\operatorname{Hom}_S(X,M)$（对每个 $S'\to S$ 及 $x\in X(S')$，有 $(m\cdot\operatorname{id}_X)(x)=m(x)\cdot x$）。而根据 3.1 的证明，每个 $m:X\to M$ 都唯一地\Term{factorization}{分解}为从 $Y=X^+$ 到 $M$ 的\Term{morphism}{态射} $h$，故 $m\cdot\operatorname{id}_X$ 就是 1.3.1 引入的\Term{automorphism}{自同构} $u_h$。于是由\Term{equivalence}{等价}的定义（见 1.3.4、1.3.5）即得\Term{lemma}{引理}。}
\end{proof}
结合前述全部结果，命题 1.3.6 给出：
\begin{SGAresult}{\Term{theorem}{定理}}{3.5}
设 $S$ 为\Term{scheme}{概形}，$\mathcal I,\mathcal J$ 为 $S$ 上两个\Term{ideal}{理想}，满足 $\mathcal I\supset\mathcal J$、$\mathcal I\cdot\mathcal J=0$；$S_0,S_J$ 为它们定义的\Term{closed-subscheme}{闭子概形}。设 $X$ 为在 $S$ 上\Term{flat}{平坦}的 $S$-\Term{scheme}{概形}（且在 $S$ 上\Term{locally-finite-presentation}{局部有限表现}或\Term{quasi-compact}{拟紧}），$X_0,X_J$ 为经\Term{base-change}{基变换}得到的\Term{scheme}{概形}。假设 $X_J$ 带有 $S_J$-\Term{group-structure}{群结构}，并以 $L_0$ 表示由以下公式定义的\Term{abelian-group}{交换群}值 $S_0$-\Term{functor}{函子}：
\[
 \operatorname{Hom}_{S_0}(T,L_0)=
 \operatorname{Hom}_{\mathcal O_T}
 \bigl(\omega^1_{X_0/S_0}\otimes_{\mathcal O_{S_0}}\mathcal O_T,
 \mathcal J\otimes_{\mathcal O_{S_0}}\mathcal O_T\bigr),
\]
其中 $X_0$ 经其\Term{adjoint-representation}{伴随表示}\Term{action}{作用}于 $L_0$。

(i) $X$ 上存在诱导 $X_J$ 上给定结构的 $S$-\Term{group-structure}{群结构}，当且仅当以下条件成立：

$(i_1)$ 存在一个诱导 $X_J$ 的\Term{group-law}{群运算}的 $S$-\Term{scheme}{概形}的\Term{morphism}{态射} $X\times_S X\to X$。

$(i_2)$ 属于 $H^3(X_0,L_0)$ 的某个\Term{obstruction-class}{障碍类}为零（该类由 $X$ 及 $X_J$ 的\Term{group-law}{群运算}\Term{canonical}{典范}定义）。

(ii) 若 (i) 的条件成立，则 $X$ 上诱导 $X_J$ 的给定\Term{group-law}{群运算}的\Term{group-law}{群运算}构成的\Term{set}{集合} $E$，是 $Z^2(X_0,L_0)$ 下的\Term{principal-homogeneous-set}{主齐性集合}；将 $E$ 按 $X$ 的那些在 $X_J$ 上诱导\Term{identity-morphism}{恒等态射}的 $S$-\Term{automorphism}{自同构}取商，所得\Term{set}{集合}是 $H^2(X_0,L_0)$ 下的\Term{principal-homogeneous-set}{主齐性集合}。
\end{SGAresult}
\begin{proof}
\footnote{编者注（69）：以下内容为新增。}事实上，每个诱导 $X_J$ 的\Term{group-law}{群运算}的 $S$-\Term{scheme}{概形}的\Term{morphism}{态射} $f:X\times_S X\to X$，由 3.2 都是 $X$ 上的\Term{admissible-law}{可容许}\Term{composition-law}{合成律}；于是由 1.3.6 (i)，存在一个结合的\Term{admissible-law}{可容许}\Term{composition-law}{合成律} $P:X\times_S X\to X$，\Term{equivalence}{等价}于某个类 $c(f)\in H^3(X_0,L_0)$ 为零，而此时根据 3.3，$P$ 为\Term{group-law}{群运算}。这证明了 (i)；再由 3.3 及 1.3.6 (ii) 得到 (ii)。
\end{proof}
\textbf{注记 3.5.1。}\footnote{编者注（70）：增加了此注记，作为 4.5.1 的对应，以引入 4.38 所用的记号 $\delta(\mu,\mu')$（或 $\delta(X,X')$）；因此也在上面的 3.5 (ii) 中增加了关于 $E$ 本身的部分。}若 $\mu,\mu'$ 为 $X$ 上诱导 $X_J$ 的给定\Term{group-law}{群运算}的\Term{group-law}{群运算}，便得到\Term{cocycle}{上闭链} $\delta(\mu,\mu')\in Z^2(X_0,L_0)$，所选\Term{signature}{符号}约定是
\SourcePage{147}
$\mu'=\delta(\mu,\mu')\cdot\mu$，即对每个 $S'\to S$ 及 $x,y\in X(S')$，有
\[
 \mu'(x,y)=\delta(\mu,\mu')(x_0,y_0)\cdot\mu(x,y).
\]\footnote{编者注（71）：我们沿用原书的\Term{signature}{符号}约定，以便在 4.38 (5) 中有等式 $\partial^1d(X,X')=\delta(X,X')$（另见编者注（54））。}
以 $\overline\delta(\mu,\mu')$ 表示 $\delta(\mu,\mu')$ 在 $H^2(X_0,L_0)$ 中的像。最后，若将带\Term{group-law}{群运算} $\mu$（相应地，$\mu'$）的 $X$ 简记为 $X$（相应地，$X'$），便用 $\delta(X,X')$ 代替 $\delta(\mu,\mu')$；$\overline\delta(X,X')$ 的记法亦同。

\textbf{注记 3.6。}设 $X_J$ 是在 $S_J$ 上\Term{smooth}{光滑}且\Term{affine}{仿射}的 $S_J$-\Term{scheme}{概形}。由 0.15，在\Term{isomorphism}{同构}意义下存在唯一在 $S$ 上\Term{smooth}{光滑}且\Term{reduced}{约化}为 $X_J$ 的 $S$-\Term{scheme}{概形} $X$。若 $X_J$ 带有 $S_J$-\Term{group-structure}{群结构}，由 0.16 知条件 $(i_1)$ 自动成立。此外，由 0.6，$L_0$ 的定义可简化，得到：
\begin{SGAresult}{推论}{3.7}
设 $S,\mathcal I,\mathcal J$ 如 3.1。设 $X_J$ 为在 $S_J$ 上\Term{smooth}{光滑}且\Term{affine}{仿射}的 $S_J$-\Term{group}{群}。

(i) 在 $S$ 上\Term{smooth}{光滑}且\Term{reduced}{约化}为 $X_J$ 的 $S$-\Term{group}{群}，按（在 $X_J$ 上诱导\Term{identity-morphism}{恒等态射}的）\Term{isomorphism}{同构}分类，所得\Term{set}{集合}或者为空，或者是\Term{group}{群}
\[
 H^2\bigl(X_0,\operatorname{Lie}(X_0/S_0)\otimes_{\mathcal O_{S_0}}\mathcal J\bigr)
\]
下的\Term{principal-homogeneous-set}{主齐性集合}。

(ii) 存在一个在 $S$ 上\Term{smooth}{光滑}且\Term{reduced}{约化}为 $X_J$ 的 $S$-\Term{group}{群}，当且仅当
\[
 H^3\bigl(X_0,\operatorname{Lie}(X_0/S_0)\otimes_{\mathcal O_{S_0}}\mathcal J\bigr)
\]
中的某个\Term{obstruction}{障碍}为零。
\end{SGAresult}
照常得到以下推论。
\begin{SGAresult}{推论}{3.8}
设 $S$ 为\Term{scheme}{概形}，$S_0$ 为一个\Term{nilpotent}{幂零}\Term{ideal}{理想} $\mathcal I$ 定义的\Term{closed-subscheme}{闭子概形}。设 $X_0$ 是在 $S$ 上\Term{smooth}{光滑}且\Term{affine}{仿射}的 $S_0$-\Term{group}{群}。

(i) 若对每个 $n\geq0$ 都有 $H^2(X_0,\operatorname{Lie}(X_0/S_0)\otimes_{\mathcal O_{S_0}}\mathcal I^{n+1}/\mathcal I^{n+2})=0$，则两个在 $S$ 上\Term{smooth}{光滑}且\Term{reduced}{约化}为 $X_0$ 的 $S$-\Term{group}{群}必\Term{isomorphism}{同构}（可取一个在 $X_0$ 上诱导\Term{identity-morphism}{恒等态射}的\Term{isomorphism}{同构}）。

(ii) 若对每个 $n\geq0$ 都有 $H^3(X_0,\operatorname{Lie}(X_0/S_0)\otimes_{\mathcal O_{S_0}}\mathcal I^{n+1}/\mathcal I^{n+2})=0$，则存在一个在 $S$ 上\Term{smooth}{光滑}且\Term{reduced}{约化}为 $X_0$ 的 $S$-\Term{group}{群}。
\end{SGAresult}
\begin{SGAresult}{推论}{3.9}
设 $S$ 为\Term{affine}{仿射}\Term{scheme}{概形}，$S_0$ 为一个\Term{nilpotent}{幂零}\Term{ideal}{理想} $\mathcal I$ 定义的\Term{closed-subscheme}{闭子概形}。假设 $\mathcal I^{n+1}/\mathcal I^{n+2}$ 在 $S_0$ 上\Term{locally-free}{局部自由}。设 $X_0$ 为在 $S_0$ 上\Term{smooth}{光滑}且\Term{affine}{仿射}的 $S_0$-\Term{group}{群}。

(i) 若 $H^2(X_0,\operatorname{Lie}(X_0/S_0))=0$，则两个在 $S$ 上\Term{smooth}{光滑}且\Term{reduced}{约化}为 $X_0$ 的 $S$-\Term{group}{群}必\Term{isomorphism}{同构}。

(ii) 若 $H^3(X_0,\operatorname{Lie}(X_0/S_0))=0$，则存在一个在 $S$ 上\Term{smooth}{光滑}且\Term{reduced}{约化}为 $X_0$ 的 $S$-\Term{group}{群}。
\end{SGAresult}
\SourcePage{148}
\begin{SGAresult}{推论}{3.10}
设 $S_0$ 为\Term{scheme}{概形}，$S=\mathrm I_{S_0}$ 为 $S_0$ 上的\Term{dual-numbers-scheme}{对偶数概形}。设 $X_0$ 为在 $S_0$ 上\Term{smooth}{光滑}的 $S_0$-\Term{group}{群}。每个在 $S$ 上\Term{smooth}{光滑}、且 $Y_0$ 与 $X_0$ $S_0$-\Term{isomorphism}{同构}的 $S$-\Term{group}{群} $Y$，都与 $X=X_0\times_{S_0}S$ $S$-\Term{isomorphism}{同构}，其充分必要条件是 $H^2(X_0,\operatorname{Lie}(X_0/S_0))=0$。\footnote{编者注（72）：第 XXIV 讲 1.13 使用了这一点。}
\end{SGAresult}
\begin{proof}
事实上，由 3.5，这些\Term{group}{群} $Y$ 按“在 $X_0$ 上诱导\Term{identity-morphism}{恒等态射}”的 $S$-\Term{group}{群}\Term{isomorphism}{同构}分类所得的\Term{set}{集合}，与 $H^2(X_0,\operatorname{Lie}(X_0/S_0))$ 之间存在\Term{bijection}{双射}；因此，按任意 $S$-\Term{group}{群}\Term{isomorphism}{同构}分类所得的\Term{set}{集合}，与
\[
 H^2(X_0,\operatorname{Lie}(X_0/S_0))/\Gamma_0
\]
之间存在\Term{bijection}{双射}，其中
\[
 \Gamma_0=\operatorname{Aut}_{S_0\text{-gr.}}(X_0)
\]
（它以显然的方式\Term{action}{作用}于 $H^2$）。结论随即得到。\footnote{编者注（73）：事实上，$\operatorname{Aut}_{S_0\text{-gr.}}(X_0)$ 通过\Term{group}{群}\Term{automorphism}{自同构}\Term{action}{作用}于\Term{abelian-group}{交换群} $H^2(X_0,\operatorname{Lie}(X_0/S_0))$，所以 $0$ 的\Term{orbit}{轨道}是单点集 $\{0\}$；因此，\Term{quotient-set}{商集}是单点集，当且仅当 $H^2(X_0,\operatorname{Lie}(X_0/S_0))=\{0\}$。}
\end{proof}
